% !TEX program = lualatex
\ifdefined\pg\else\IfFileExists{tex/avulso.tex}{% Compilação AVULSA de um fragmento (volN/.../<stem>.tex) — p. ex. pelo botão do editor,
% de dentro da pasta do nó.  O fragmento começa com
%
%   \ifdefined\pg\else\IfFileExists{tex/avulso.tex}{% Compilação AVULSA de um fragmento (volN/.../<stem>.tex) — p. ex. pelo botão do editor,
% de dentro da pasta do nó.  O fragmento começa com
%
%   \ifdefined\pg\else\IfFileExists{tex/avulso.tex}{% Compilação AVULSA de um fragmento (volN/.../<stem>.tex) — p. ex. pelo botão do editor,
% de dentro da pasta do nó.  O fragmento começa com
%
%   \ifdefined\pg\else\IfFileExists{tex/avulso.tex}{\input{tex/avulso}}{\input{../../tex/avulso}}\fi
%
% e termina com  \csname fimavulso\endcsname .  Dentro do livro ou do _preview, \pg já existe
% e nada disto roda; \fimavulso fica indefinido e o \csname vira \relax.
%
% Motor: o projeto é lualatex (fontspec/polyglossia).  Se o editor chamar pdflatex, cai
% num preâmbulo equivalente com inputenc/babel — serve para conferir, não para o PDF final.
\documentclass[11pt,a5paper]{book}
\ifdefined\directlua
  \usepackage{fontspec}
  \usepackage{polyglossia}
  \setmainlanguage{portuguese}
\else\ifdefined\XeTeXrevision
  \usepackage{fontspec}
  \usepackage{polyglossia}
  \setmainlanguage{portuguese}
\else
  \usepackage[utf8]{inputenc}
  \usepackage[T1]{fontenc}
  \usepackage{lmodern}
  \usepackage[portuguese]{babel}
  \DeclareUnicodeCharacter{2192}{\ensuremath{\rightarrow}}
\fi\fi
% raiz do repo: compilado da raiz ou de volN/<parte>/ (dois níveis abaixo)
\IfFileExists{tex/preambulo.tex}{\input{tex/preambulo}}{\input{../../tex/preambulo}}
\def\fimavulso{\end{document}}
\begin{document}
}{% Compilação AVULSA de um fragmento (volN/.../<stem>.tex) — p. ex. pelo botão do editor,
% de dentro da pasta do nó.  O fragmento começa com
%
%   \ifdefined\pg\else\IfFileExists{tex/avulso.tex}{\input{tex/avulso}}{\input{../../tex/avulso}}\fi
%
% e termina com  \csname fimavulso\endcsname .  Dentro do livro ou do _preview, \pg já existe
% e nada disto roda; \fimavulso fica indefinido e o \csname vira \relax.
%
% Motor: o projeto é lualatex (fontspec/polyglossia).  Se o editor chamar pdflatex, cai
% num preâmbulo equivalente com inputenc/babel — serve para conferir, não para o PDF final.
\documentclass[11pt,a5paper]{book}
\ifdefined\directlua
  \usepackage{fontspec}
  \usepackage{polyglossia}
  \setmainlanguage{portuguese}
\else\ifdefined\XeTeXrevision
  \usepackage{fontspec}
  \usepackage{polyglossia}
  \setmainlanguage{portuguese}
\else
  \usepackage[utf8]{inputenc}
  \usepackage[T1]{fontenc}
  \usepackage{lmodern}
  \usepackage[portuguese]{babel}
  \DeclareUnicodeCharacter{2192}{\ensuremath{\rightarrow}}
\fi\fi
% raiz do repo: compilado da raiz ou de volN/<parte>/ (dois níveis abaixo)
\IfFileExists{tex/preambulo.tex}{\input{tex/preambulo}}{\input{../../tex/preambulo}}
\def\fimavulso{\end{document}}
\begin{document}
}\fi
%
% e termina com  \csname fimavulso\endcsname .  Dentro do livro ou do _preview, \pg já existe
% e nada disto roda; \fimavulso fica indefinido e o \csname vira \relax.
%
% Motor: o projeto é lualatex (fontspec/polyglossia).  Se o editor chamar pdflatex, cai
% num preâmbulo equivalente com inputenc/babel — serve para conferir, não para o PDF final.
\documentclass[11pt,a5paper]{book}
\ifdefined\directlua
  \usepackage{fontspec}
  \usepackage{polyglossia}
  \setmainlanguage{portuguese}
\else\ifdefined\XeTeXrevision
  \usepackage{fontspec}
  \usepackage{polyglossia}
  \setmainlanguage{portuguese}
\else
  \usepackage[utf8]{inputenc}
  \usepackage[T1]{fontenc}
  \usepackage{lmodern}
  \usepackage[portuguese]{babel}
  \DeclareUnicodeCharacter{2192}{\ensuremath{\rightarrow}}
\fi\fi
% raiz do repo: compilado da raiz ou de volN/<parte>/ (dois níveis abaixo)
\IfFileExists{tex/preambulo.tex}{% Preâmbulo comum ao livro e à pré-visualização de fragmentos (tex/_preview.tex).

% Mancha em A5.  Sem isto vale o padrão da classe book (~105 mm de largura e margem
% externa de ~25 mm em frente e verso), largo demais para o formato.  Aqui: ~118 mm de
% largura, margem interna um pouco maior que a externa (encadernação).
\usepackage[a5paper,inner=16mm,outer=14mm,top=14mm,bottom=16mm,
            includehead,headsep=5mm,footskip=8mm]{geometry}

% Títulos de secção sem justificar (evita linhas estouradas em \Large bold, p. ex.
% «Demonstrabilidade») e um pouco de folga para enunciados de tese em negrito.
\usepackage{sectsty}
\allsectionsfont{\raggedright}
\setlength{\emergencystretch}{1.5em}

% Marca da página impressa do original: \pg{286} imprime [286] no fluxo do texto.
\newcommand{\pg}[1]{\,{\footnotesize\textbf{[#1]}}\,}

% Virada de página do original DENTRO de uma nota (\pgn{422}); discreta, para não quebrar a nota.
\newcommand{\pgn}[1]{{\scriptsize\textbf{[#1]}}}

% Asterisco do original (divisão de assunto, sem título): \asterisco centrado em linha própria,
% também válido no meio do parágrafo ou dentro de nota (\footnote{... \asterisco ...}).
\newcommand{\asterisco}{\par\medskip\centerline{\footnotesize$\ast$}\medskip\noindent}

% Gravura recortada do fac-símile (<figure> do TEI; tools/tei2latex.py): fora de float, no ponto
% exato do original.  \gravura[detalhes]{arquivo.png}{fração da largura do texto}{legenda}:
% legenda principal centrada ("FIGURA VIII. — Côrte…"); os detalhes opcionais (<p> do <figure>:
% "A, hemispherio…; — B, …") vão abaixo, em corpo menor e justificados.  Ambos podem ser vazios.
\usepackage{graphicx}
\newcommand{\gravura}[4][]{\par\medskip\begingroup\centering
  \includegraphics[width=#3\linewidth,height=0.75\textheight,keepaspectratio]{#2}\par
  \if\relax\detokenize{#4}\relax\else\smallskip{\small #4\par}\fi
  \if\relax\detokenize{#1}\relax\else\smallskip{\footnotesize\leftskip=1em\rightskip=1em
    \parfillskip=0pt plus 1fil\noindent #1\par}\fi\endgroup\medskip\noindent}

% Citações latinas de S. Tomás (o OCR erra muito aqui: sempre conferir na imagem).
\newenvironment{latim}{\begin{quote}\itshape}{\end{quote}}

% --- Estrutura do autor (promovido de "Artigo III.tex") -----------------------
% enunciado de tese numerada que corre com o texto (segue-se travessão)
\newcommand{\tese}[2]{\bigskip\noindent\textbf{#1.\quad #2}\;\textemdash\;}
% enunciado de tese numerada em linha própria (seguem-se as subdivisões a), b)...)
\newcommand{\teseB}[2]{\bigskip\noindent\textbf{#1.\quad #2}\par}
% subdivisão a), b), c)... com rubrica em itálico
\newcommand{\sub}[2]{\medskip\noindent#1)\quad\emph{#2},\;\textemdash\;}
% subdivisão a), b), c)... sem rubrica
\newcommand{\subp}[1]{\medskip\noindent#1)\quad}

% --- Notas: no pé (padrão) ou correndo no corpo do texto ---------------------
% Passar  \def\notascorridas{}  no -usepretex faz cada nota correr no corpo, em
% corpo menor, num bloco recuado que interrompe a leitura até acabar.  Ex.:
%   latexmk -lualatex -usepretex='\def\fragmento{vol2/.../01_cap_y.tex}\def\notascorridas{}' tex/_preview.tex
\newif\ifnotascorridas
\ifdefined\notascorridas \notascorridastrue \fi
\ifnotascorridas
  \renewcommand\footnote[1]{%
    \stepcounter{footnote}\textsuperscript{\thefootnote}% marca no ponto do texto
    \par
    {\footnotesize
     \leftskip=2em \rightskip=2em \parindent=1em
     \addvspace{\medskipamount}%
     \noindent\textit{Nota \thefootnote.}\enspace #1\par
     \addvspace{\medskipamount}}%
  }
\fi

% --- Sem notas: \def\semnotas{} no -usepretex suprime todas as notas ---------
% (nem a chamada no texto nem o conteúdo).  Usado por  ./compila.sh --sem-notas .
\ifdefined\semnotas
  \renewcommand\footnote[2][]{}%
  \renewcommand\footnotemark[1][]{}%
  \renewcommand\footnotetext[2][]{}%
\fi

% peças pré-textuais (nó pretextual; p@type em web/obra/esquema.py)
\newcommand{\preTituloObra}[1]{\par\vspace{1em}{\centering\Large\bfseries\scshape #1\par}\vspace{1em}}
\newcommand{\preCentro}[1]{\par{\centering #1\par}}
\newcommand{\preEnderecamento}[1]{\par\vspace{1em}{\centering\small\scshape #1\par}}
\newcommand{\preSaudacao}[1]{\par\medskip\noindent #1\par}
\newcommand{\preDataLocal}[1]{\par\medskip\noindent\emph{#1}\par}
\newcommand{\preFecho}[1]{\par\smallskip{\raggedleft\scshape #1\par}}
}{% Preâmbulo comum ao livro e à pré-visualização de fragmentos (tex/_preview.tex).

% Mancha em A5.  Sem isto vale o padrão da classe book (~105 mm de largura e margem
% externa de ~25 mm em frente e verso), largo demais para o formato.  Aqui: ~118 mm de
% largura, margem interna um pouco maior que a externa (encadernação).
\usepackage[a5paper,inner=16mm,outer=14mm,top=14mm,bottom=16mm,
            includehead,headsep=5mm,footskip=8mm]{geometry}

% Títulos de secção sem justificar (evita linhas estouradas em \Large bold, p. ex.
% «Demonstrabilidade») e um pouco de folga para enunciados de tese em negrito.
\usepackage{sectsty}
\allsectionsfont{\raggedright}
\setlength{\emergencystretch}{1.5em}

% Marca da página impressa do original: \pg{286} imprime [286] no fluxo do texto.
\newcommand{\pg}[1]{\,{\footnotesize\textbf{[#1]}}\,}

% Virada de página do original DENTRO de uma nota (\pgn{422}); discreta, para não quebrar a nota.
\newcommand{\pgn}[1]{{\scriptsize\textbf{[#1]}}}

% Asterisco do original (divisão de assunto, sem título): \asterisco centrado em linha própria,
% também válido no meio do parágrafo ou dentro de nota (\footnote{... \asterisco ...}).
\newcommand{\asterisco}{\par\medskip\centerline{\footnotesize$\ast$}\medskip\noindent}

% Gravura recortada do fac-símile (<figure> do TEI; tools/tei2latex.py): fora de float, no ponto
% exato do original.  \gravura[detalhes]{arquivo.png}{fração da largura do texto}{legenda}:
% legenda principal centrada ("FIGURA VIII. — Côrte…"); os detalhes opcionais (<p> do <figure>:
% "A, hemispherio…; — B, …") vão abaixo, em corpo menor e justificados.  Ambos podem ser vazios.
\usepackage{graphicx}
\newcommand{\gravura}[4][]{\par\medskip\begingroup\centering
  \includegraphics[width=#3\linewidth,height=0.75\textheight,keepaspectratio]{#2}\par
  \if\relax\detokenize{#4}\relax\else\smallskip{\small #4\par}\fi
  \if\relax\detokenize{#1}\relax\else\smallskip{\footnotesize\leftskip=1em\rightskip=1em
    \parfillskip=0pt plus 1fil\noindent #1\par}\fi\endgroup\medskip\noindent}

% Citações latinas de S. Tomás (o OCR erra muito aqui: sempre conferir na imagem).
\newenvironment{latim}{\begin{quote}\itshape}{\end{quote}}

% --- Estrutura do autor (promovido de "Artigo III.tex") -----------------------
% enunciado de tese numerada que corre com o texto (segue-se travessão)
\newcommand{\tese}[2]{\bigskip\noindent\textbf{#1.\quad #2}\;\textemdash\;}
% enunciado de tese numerada em linha própria (seguem-se as subdivisões a), b)...)
\newcommand{\teseB}[2]{\bigskip\noindent\textbf{#1.\quad #2}\par}
% subdivisão a), b), c)... com rubrica em itálico
\newcommand{\sub}[2]{\medskip\noindent#1)\quad\emph{#2},\;\textemdash\;}
% subdivisão a), b), c)... sem rubrica
\newcommand{\subp}[1]{\medskip\noindent#1)\quad}

% --- Notas: no pé (padrão) ou correndo no corpo do texto ---------------------
% Passar  \def\notascorridas{}  no -usepretex faz cada nota correr no corpo, em
% corpo menor, num bloco recuado que interrompe a leitura até acabar.  Ex.:
%   latexmk -lualatex -usepretex='\def\fragmento{vol2/.../01_cap_y.tex}\def\notascorridas{}' tex/_preview.tex
\newif\ifnotascorridas
\ifdefined\notascorridas \notascorridastrue \fi
\ifnotascorridas
  \renewcommand\footnote[1]{%
    \stepcounter{footnote}\textsuperscript{\thefootnote}% marca no ponto do texto
    \par
    {\footnotesize
     \leftskip=2em \rightskip=2em \parindent=1em
     \addvspace{\medskipamount}%
     \noindent\textit{Nota \thefootnote.}\enspace #1\par
     \addvspace{\medskipamount}}%
  }
\fi

% --- Sem notas: \def\semnotas{} no -usepretex suprime todas as notas ---------
% (nem a chamada no texto nem o conteúdo).  Usado por  ./compila.sh --sem-notas .
\ifdefined\semnotas
  \renewcommand\footnote[2][]{}%
  \renewcommand\footnotemark[1][]{}%
  \renewcommand\footnotetext[2][]{}%
\fi

% peças pré-textuais (nó pretextual; p@type em web/obra/esquema.py)
\newcommand{\preTituloObra}[1]{\par\vspace{1em}{\centering\Large\bfseries\scshape #1\par}\vspace{1em}}
\newcommand{\preCentro}[1]{\par{\centering #1\par}}
\newcommand{\preEnderecamento}[1]{\par\vspace{1em}{\centering\small\scshape #1\par}}
\newcommand{\preSaudacao}[1]{\par\medskip\noindent #1\par}
\newcommand{\preDataLocal}[1]{\par\medskip\noindent\emph{#1}\par}
\newcommand{\preFecho}[1]{\par\smallskip{\raggedleft\scshape #1\par}}
}
\def\fimavulso{\end{document}}
\begin{document}
}{% Compilação AVULSA de um fragmento (volN/.../<stem>.tex) — p. ex. pelo botão do editor,
% de dentro da pasta do nó.  O fragmento começa com
%
%   \ifdefined\pg\else\IfFileExists{tex/avulso.tex}{% Compilação AVULSA de um fragmento (volN/.../<stem>.tex) — p. ex. pelo botão do editor,
% de dentro da pasta do nó.  O fragmento começa com
%
%   \ifdefined\pg\else\IfFileExists{tex/avulso.tex}{\input{tex/avulso}}{\input{../../tex/avulso}}\fi
%
% e termina com  \csname fimavulso\endcsname .  Dentro do livro ou do _preview, \pg já existe
% e nada disto roda; \fimavulso fica indefinido e o \csname vira \relax.
%
% Motor: o projeto é lualatex (fontspec/polyglossia).  Se o editor chamar pdflatex, cai
% num preâmbulo equivalente com inputenc/babel — serve para conferir, não para o PDF final.
\documentclass[11pt,a5paper]{book}
\ifdefined\directlua
  \usepackage{fontspec}
  \usepackage{polyglossia}
  \setmainlanguage{portuguese}
\else\ifdefined\XeTeXrevision
  \usepackage{fontspec}
  \usepackage{polyglossia}
  \setmainlanguage{portuguese}
\else
  \usepackage[utf8]{inputenc}
  \usepackage[T1]{fontenc}
  \usepackage{lmodern}
  \usepackage[portuguese]{babel}
  \DeclareUnicodeCharacter{2192}{\ensuremath{\rightarrow}}
\fi\fi
% raiz do repo: compilado da raiz ou de volN/<parte>/ (dois níveis abaixo)
\IfFileExists{tex/preambulo.tex}{\input{tex/preambulo}}{\input{../../tex/preambulo}}
\def\fimavulso{\end{document}}
\begin{document}
}{% Compilação AVULSA de um fragmento (volN/.../<stem>.tex) — p. ex. pelo botão do editor,
% de dentro da pasta do nó.  O fragmento começa com
%
%   \ifdefined\pg\else\IfFileExists{tex/avulso.tex}{\input{tex/avulso}}{\input{../../tex/avulso}}\fi
%
% e termina com  \csname fimavulso\endcsname .  Dentro do livro ou do _preview, \pg já existe
% e nada disto roda; \fimavulso fica indefinido e o \csname vira \relax.
%
% Motor: o projeto é lualatex (fontspec/polyglossia).  Se o editor chamar pdflatex, cai
% num preâmbulo equivalente com inputenc/babel — serve para conferir, não para o PDF final.
\documentclass[11pt,a5paper]{book}
\ifdefined\directlua
  \usepackage{fontspec}
  \usepackage{polyglossia}
  \setmainlanguage{portuguese}
\else\ifdefined\XeTeXrevision
  \usepackage{fontspec}
  \usepackage{polyglossia}
  \setmainlanguage{portuguese}
\else
  \usepackage[utf8]{inputenc}
  \usepackage[T1]{fontenc}
  \usepackage{lmodern}
  \usepackage[portuguese]{babel}
  \DeclareUnicodeCharacter{2192}{\ensuremath{\rightarrow}}
\fi\fi
% raiz do repo: compilado da raiz ou de volN/<parte>/ (dois níveis abaixo)
\IfFileExists{tex/preambulo.tex}{\input{tex/preambulo}}{\input{../../tex/preambulo}}
\def\fimavulso{\end{document}}
\begin{document}
}\fi
%
% e termina com  \csname fimavulso\endcsname .  Dentro do livro ou do _preview, \pg já existe
% e nada disto roda; \fimavulso fica indefinido e o \csname vira \relax.
%
% Motor: o projeto é lualatex (fontspec/polyglossia).  Se o editor chamar pdflatex, cai
% num preâmbulo equivalente com inputenc/babel — serve para conferir, não para o PDF final.
\documentclass[11pt,a5paper]{book}
\ifdefined\directlua
  \usepackage{fontspec}
  \usepackage{polyglossia}
  \setmainlanguage{portuguese}
\else\ifdefined\XeTeXrevision
  \usepackage{fontspec}
  \usepackage{polyglossia}
  \setmainlanguage{portuguese}
\else
  \usepackage[utf8]{inputenc}
  \usepackage[T1]{fontenc}
  \usepackage{lmodern}
  \usepackage[portuguese]{babel}
  \DeclareUnicodeCharacter{2192}{\ensuremath{\rightarrow}}
\fi\fi
% raiz do repo: compilado da raiz ou de volN/<parte>/ (dois níveis abaixo)
\IfFileExists{tex/preambulo.tex}{% Preâmbulo comum ao livro e à pré-visualização de fragmentos (tex/_preview.tex).

% Mancha em A5.  Sem isto vale o padrão da classe book (~105 mm de largura e margem
% externa de ~25 mm em frente e verso), largo demais para o formato.  Aqui: ~118 mm de
% largura, margem interna um pouco maior que a externa (encadernação).
\usepackage[a5paper,inner=16mm,outer=14mm,top=14mm,bottom=16mm,
            includehead,headsep=5mm,footskip=8mm]{geometry}

% Títulos de secção sem justificar (evita linhas estouradas em \Large bold, p. ex.
% «Demonstrabilidade») e um pouco de folga para enunciados de tese em negrito.
\usepackage{sectsty}
\allsectionsfont{\raggedright}
\setlength{\emergencystretch}{1.5em}

% Marca da página impressa do original: \pg{286} imprime [286] no fluxo do texto.
\newcommand{\pg}[1]{\,{\footnotesize\textbf{[#1]}}\,}

% Virada de página do original DENTRO de uma nota (\pgn{422}); discreta, para não quebrar a nota.
\newcommand{\pgn}[1]{{\scriptsize\textbf{[#1]}}}

% Asterisco do original (divisão de assunto, sem título): \asterisco centrado em linha própria,
% também válido no meio do parágrafo ou dentro de nota (\footnote{... \asterisco ...}).
\newcommand{\asterisco}{\par\medskip\centerline{\footnotesize$\ast$}\medskip\noindent}

% Gravura recortada do fac-símile (<figure> do TEI; tools/tei2latex.py): fora de float, no ponto
% exato do original.  \gravura[detalhes]{arquivo.png}{fração da largura do texto}{legenda}:
% legenda principal centrada ("FIGURA VIII. — Côrte…"); os detalhes opcionais (<p> do <figure>:
% "A, hemispherio…; — B, …") vão abaixo, em corpo menor e justificados.  Ambos podem ser vazios.
\usepackage{graphicx}
\newcommand{\gravura}[4][]{\par\medskip\begingroup\centering
  \includegraphics[width=#3\linewidth,height=0.75\textheight,keepaspectratio]{#2}\par
  \if\relax\detokenize{#4}\relax\else\smallskip{\small #4\par}\fi
  \if\relax\detokenize{#1}\relax\else\smallskip{\footnotesize\leftskip=1em\rightskip=1em
    \parfillskip=0pt plus 1fil\noindent #1\par}\fi\endgroup\medskip\noindent}

% Citações latinas de S. Tomás (o OCR erra muito aqui: sempre conferir na imagem).
\newenvironment{latim}{\begin{quote}\itshape}{\end{quote}}

% --- Estrutura do autor (promovido de "Artigo III.tex") -----------------------
% enunciado de tese numerada que corre com o texto (segue-se travessão)
\newcommand{\tese}[2]{\bigskip\noindent\textbf{#1.\quad #2}\;\textemdash\;}
% enunciado de tese numerada em linha própria (seguem-se as subdivisões a), b)...)
\newcommand{\teseB}[2]{\bigskip\noindent\textbf{#1.\quad #2}\par}
% subdivisão a), b), c)... com rubrica em itálico
\newcommand{\sub}[2]{\medskip\noindent#1)\quad\emph{#2},\;\textemdash\;}
% subdivisão a), b), c)... sem rubrica
\newcommand{\subp}[1]{\medskip\noindent#1)\quad}

% --- Notas: no pé (padrão) ou correndo no corpo do texto ---------------------
% Passar  \def\notascorridas{}  no -usepretex faz cada nota correr no corpo, em
% corpo menor, num bloco recuado que interrompe a leitura até acabar.  Ex.:
%   latexmk -lualatex -usepretex='\def\fragmento{vol2/.../01_cap_y.tex}\def\notascorridas{}' tex/_preview.tex
\newif\ifnotascorridas
\ifdefined\notascorridas \notascorridastrue \fi
\ifnotascorridas
  \renewcommand\footnote[1]{%
    \stepcounter{footnote}\textsuperscript{\thefootnote}% marca no ponto do texto
    \par
    {\footnotesize
     \leftskip=2em \rightskip=2em \parindent=1em
     \addvspace{\medskipamount}%
     \noindent\textit{Nota \thefootnote.}\enspace #1\par
     \addvspace{\medskipamount}}%
  }
\fi

% --- Sem notas: \def\semnotas{} no -usepretex suprime todas as notas ---------
% (nem a chamada no texto nem o conteúdo).  Usado por  ./compila.sh --sem-notas .
\ifdefined\semnotas
  \renewcommand\footnote[2][]{}%
  \renewcommand\footnotemark[1][]{}%
  \renewcommand\footnotetext[2][]{}%
\fi

% peças pré-textuais (nó pretextual; p@type em web/obra/esquema.py)
\newcommand{\preTituloObra}[1]{\par\vspace{1em}{\centering\Large\bfseries\scshape #1\par}\vspace{1em}}
\newcommand{\preCentro}[1]{\par{\centering #1\par}}
\newcommand{\preEnderecamento}[1]{\par\vspace{1em}{\centering\small\scshape #1\par}}
\newcommand{\preSaudacao}[1]{\par\medskip\noindent #1\par}
\newcommand{\preDataLocal}[1]{\par\medskip\noindent\emph{#1}\par}
\newcommand{\preFecho}[1]{\par\smallskip{\raggedleft\scshape #1\par}}
}{% Preâmbulo comum ao livro e à pré-visualização de fragmentos (tex/_preview.tex).

% Mancha em A5.  Sem isto vale o padrão da classe book (~105 mm de largura e margem
% externa de ~25 mm em frente e verso), largo demais para o formato.  Aqui: ~118 mm de
% largura, margem interna um pouco maior que a externa (encadernação).
\usepackage[a5paper,inner=16mm,outer=14mm,top=14mm,bottom=16mm,
            includehead,headsep=5mm,footskip=8mm]{geometry}

% Títulos de secção sem justificar (evita linhas estouradas em \Large bold, p. ex.
% «Demonstrabilidade») e um pouco de folga para enunciados de tese em negrito.
\usepackage{sectsty}
\allsectionsfont{\raggedright}
\setlength{\emergencystretch}{1.5em}

% Marca da página impressa do original: \pg{286} imprime [286] no fluxo do texto.
\newcommand{\pg}[1]{\,{\footnotesize\textbf{[#1]}}\,}

% Virada de página do original DENTRO de uma nota (\pgn{422}); discreta, para não quebrar a nota.
\newcommand{\pgn}[1]{{\scriptsize\textbf{[#1]}}}

% Asterisco do original (divisão de assunto, sem título): \asterisco centrado em linha própria,
% também válido no meio do parágrafo ou dentro de nota (\footnote{... \asterisco ...}).
\newcommand{\asterisco}{\par\medskip\centerline{\footnotesize$\ast$}\medskip\noindent}

% Gravura recortada do fac-símile (<figure> do TEI; tools/tei2latex.py): fora de float, no ponto
% exato do original.  \gravura[detalhes]{arquivo.png}{fração da largura do texto}{legenda}:
% legenda principal centrada ("FIGURA VIII. — Côrte…"); os detalhes opcionais (<p> do <figure>:
% "A, hemispherio…; — B, …") vão abaixo, em corpo menor e justificados.  Ambos podem ser vazios.
\usepackage{graphicx}
\newcommand{\gravura}[4][]{\par\medskip\begingroup\centering
  \includegraphics[width=#3\linewidth,height=0.75\textheight,keepaspectratio]{#2}\par
  \if\relax\detokenize{#4}\relax\else\smallskip{\small #4\par}\fi
  \if\relax\detokenize{#1}\relax\else\smallskip{\footnotesize\leftskip=1em\rightskip=1em
    \parfillskip=0pt plus 1fil\noindent #1\par}\fi\endgroup\medskip\noindent}

% Citações latinas de S. Tomás (o OCR erra muito aqui: sempre conferir na imagem).
\newenvironment{latim}{\begin{quote}\itshape}{\end{quote}}

% --- Estrutura do autor (promovido de "Artigo III.tex") -----------------------
% enunciado de tese numerada que corre com o texto (segue-se travessão)
\newcommand{\tese}[2]{\bigskip\noindent\textbf{#1.\quad #2}\;\textemdash\;}
% enunciado de tese numerada em linha própria (seguem-se as subdivisões a), b)...)
\newcommand{\teseB}[2]{\bigskip\noindent\textbf{#1.\quad #2}\par}
% subdivisão a), b), c)... com rubrica em itálico
\newcommand{\sub}[2]{\medskip\noindent#1)\quad\emph{#2},\;\textemdash\;}
% subdivisão a), b), c)... sem rubrica
\newcommand{\subp}[1]{\medskip\noindent#1)\quad}

% --- Notas: no pé (padrão) ou correndo no corpo do texto ---------------------
% Passar  \def\notascorridas{}  no -usepretex faz cada nota correr no corpo, em
% corpo menor, num bloco recuado que interrompe a leitura até acabar.  Ex.:
%   latexmk -lualatex -usepretex='\def\fragmento{vol2/.../01_cap_y.tex}\def\notascorridas{}' tex/_preview.tex
\newif\ifnotascorridas
\ifdefined\notascorridas \notascorridastrue \fi
\ifnotascorridas
  \renewcommand\footnote[1]{%
    \stepcounter{footnote}\textsuperscript{\thefootnote}% marca no ponto do texto
    \par
    {\footnotesize
     \leftskip=2em \rightskip=2em \parindent=1em
     \addvspace{\medskipamount}%
     \noindent\textit{Nota \thefootnote.}\enspace #1\par
     \addvspace{\medskipamount}}%
  }
\fi

% --- Sem notas: \def\semnotas{} no -usepretex suprime todas as notas ---------
% (nem a chamada no texto nem o conteúdo).  Usado por  ./compila.sh --sem-notas .
\ifdefined\semnotas
  \renewcommand\footnote[2][]{}%
  \renewcommand\footnotemark[1][]{}%
  \renewcommand\footnotetext[2][]{}%
\fi

% peças pré-textuais (nó pretextual; p@type em web/obra/esquema.py)
\newcommand{\preTituloObra}[1]{\par\vspace{1em}{\centering\Large\bfseries\scshape #1\par}\vspace{1em}}
\newcommand{\preCentro}[1]{\par{\centering #1\par}}
\newcommand{\preEnderecamento}[1]{\par\vspace{1em}{\centering\small\scshape #1\par}}
\newcommand{\preSaudacao}[1]{\par\medskip\noindent #1\par}
\newcommand{\preDataLocal}[1]{\par\medskip\noindent\emph{#1}\par}
\newcommand{\preFecho}[1]{\par\smallskip{\raggedleft\scshape #1\par}}
}
\def\fimavulso{\end{document}}
\begin{document}
}\fi
%
% e termina com  \csname fimavulso\endcsname .  Dentro do livro ou do _preview, \pg já existe
% e nada disto roda; \fimavulso fica indefinido e o \csname vira \relax.
%
% Motor: o projeto é lualatex (fontspec/polyglossia).  Se o editor chamar pdflatex, cai
% num preâmbulo equivalente com inputenc/babel — serve para conferir, não para o PDF final.
\documentclass[11pt,a5paper]{book}
\ifdefined\directlua
  \usepackage{fontspec}
  \usepackage{polyglossia}
  \setmainlanguage{portuguese}
\else\ifdefined\XeTeXrevision
  \usepackage{fontspec}
  \usepackage{polyglossia}
  \setmainlanguage{portuguese}
\else
  \usepackage[utf8]{inputenc}
  \usepackage[T1]{fontenc}
  \usepackage{lmodern}
  \usepackage[portuguese]{babel}
  \DeclareUnicodeCharacter{2192}{\ensuremath{\rightarrow}}
\fi\fi
% raiz do repo: compilado da raiz ou de volN/<parte>/ (dois níveis abaixo)
\IfFileExists{tex/preambulo.tex}{% Preâmbulo comum ao livro e à pré-visualização de fragmentos (tex/_preview.tex).

% Mancha em A5.  Sem isto vale o padrão da classe book (~105 mm de largura e margem
% externa de ~25 mm em frente e verso), largo demais para o formato.  Aqui: ~118 mm de
% largura, margem interna um pouco maior que a externa (encadernação).
\usepackage[a5paper,inner=16mm,outer=14mm,top=14mm,bottom=16mm,
            includehead,headsep=5mm,footskip=8mm]{geometry}

% Títulos de secção sem justificar (evita linhas estouradas em \Large bold, p. ex.
% «Demonstrabilidade») e um pouco de folga para enunciados de tese em negrito.
\usepackage{sectsty}
\allsectionsfont{\raggedright}
\setlength{\emergencystretch}{1.5em}

% Marca da página impressa do original: \pg{286} imprime [286] no fluxo do texto.
\newcommand{\pg}[1]{\,{\footnotesize\textbf{[#1]}}\,}

% Virada de página do original DENTRO de uma nota (\pgn{422}); discreta, para não quebrar a nota.
\newcommand{\pgn}[1]{{\scriptsize\textbf{[#1]}}}

% Asterisco do original (divisão de assunto, sem título): \asterisco centrado em linha própria,
% também válido no meio do parágrafo ou dentro de nota (\footnote{... \asterisco ...}).
\newcommand{\asterisco}{\par\medskip\centerline{\footnotesize$\ast$}\medskip\noindent}

% Gravura recortada do fac-símile (<figure> do TEI; tools/tei2latex.py): fora de float, no ponto
% exato do original.  \gravura[detalhes]{arquivo.png}{fração da largura do texto}{legenda}:
% legenda principal centrada ("FIGURA VIII. — Côrte…"); os detalhes opcionais (<p> do <figure>:
% "A, hemispherio…; — B, …") vão abaixo, em corpo menor e justificados.  Ambos podem ser vazios.
\usepackage{graphicx}
\newcommand{\gravura}[4][]{\par\medskip\begingroup\centering
  \includegraphics[width=#3\linewidth,height=0.75\textheight,keepaspectratio]{#2}\par
  \if\relax\detokenize{#4}\relax\else\smallskip{\small #4\par}\fi
  \if\relax\detokenize{#1}\relax\else\smallskip{\footnotesize\leftskip=1em\rightskip=1em
    \parfillskip=0pt plus 1fil\noindent #1\par}\fi\endgroup\medskip\noindent}

% Citações latinas de S. Tomás (o OCR erra muito aqui: sempre conferir na imagem).
\newenvironment{latim}{\begin{quote}\itshape}{\end{quote}}

% --- Estrutura do autor (promovido de "Artigo III.tex") -----------------------
% enunciado de tese numerada que corre com o texto (segue-se travessão)
\newcommand{\tese}[2]{\bigskip\noindent\textbf{#1.\quad #2}\;\textemdash\;}
% enunciado de tese numerada em linha própria (seguem-se as subdivisões a), b)...)
\newcommand{\teseB}[2]{\bigskip\noindent\textbf{#1.\quad #2}\par}
% subdivisão a), b), c)... com rubrica em itálico
\newcommand{\sub}[2]{\medskip\noindent#1)\quad\emph{#2},\;\textemdash\;}
% subdivisão a), b), c)... sem rubrica
\newcommand{\subp}[1]{\medskip\noindent#1)\quad}

% --- Notas: no pé (padrão) ou correndo no corpo do texto ---------------------
% Passar  \def\notascorridas{}  no -usepretex faz cada nota correr no corpo, em
% corpo menor, num bloco recuado que interrompe a leitura até acabar.  Ex.:
%   latexmk -lualatex -usepretex='\def\fragmento{vol2/.../01_cap_y.tex}\def\notascorridas{}' tex/_preview.tex
\newif\ifnotascorridas
\ifdefined\notascorridas \notascorridastrue \fi
\ifnotascorridas
  \renewcommand\footnote[1]{%
    \stepcounter{footnote}\textsuperscript{\thefootnote}% marca no ponto do texto
    \par
    {\footnotesize
     \leftskip=2em \rightskip=2em \parindent=1em
     \addvspace{\medskipamount}%
     \noindent\textit{Nota \thefootnote.}\enspace #1\par
     \addvspace{\medskipamount}}%
  }
\fi

% --- Sem notas: \def\semnotas{} no -usepretex suprime todas as notas ---------
% (nem a chamada no texto nem o conteúdo).  Usado por  ./compila.sh --sem-notas .
\ifdefined\semnotas
  \renewcommand\footnote[2][]{}%
  \renewcommand\footnotemark[1][]{}%
  \renewcommand\footnotetext[2][]{}%
\fi

% peças pré-textuais (nó pretextual; p@type em web/obra/esquema.py)
\newcommand{\preTituloObra}[1]{\par\vspace{1em}{\centering\Large\bfseries\scshape #1\par}\vspace{1em}}
\newcommand{\preCentro}[1]{\par{\centering #1\par}}
\newcommand{\preEnderecamento}[1]{\par\vspace{1em}{\centering\small\scshape #1\par}}
\newcommand{\preSaudacao}[1]{\par\medskip\noindent #1\par}
\newcommand{\preDataLocal}[1]{\par\medskip\noindent\emph{#1}\par}
\newcommand{\preFecho}[1]{\par\smallskip{\raggedleft\scshape #1\par}}
}{% Preâmbulo comum ao livro e à pré-visualização de fragmentos (tex/_preview.tex).

% Mancha em A5.  Sem isto vale o padrão da classe book (~105 mm de largura e margem
% externa de ~25 mm em frente e verso), largo demais para o formato.  Aqui: ~118 mm de
% largura, margem interna um pouco maior que a externa (encadernação).
\usepackage[a5paper,inner=16mm,outer=14mm,top=14mm,bottom=16mm,
            includehead,headsep=5mm,footskip=8mm]{geometry}

% Títulos de secção sem justificar (evita linhas estouradas em \Large bold, p. ex.
% «Demonstrabilidade») e um pouco de folga para enunciados de tese em negrito.
\usepackage{sectsty}
\allsectionsfont{\raggedright}
\setlength{\emergencystretch}{1.5em}

% Marca da página impressa do original: \pg{286} imprime [286] no fluxo do texto.
\newcommand{\pg}[1]{\,{\footnotesize\textbf{[#1]}}\,}

% Virada de página do original DENTRO de uma nota (\pgn{422}); discreta, para não quebrar a nota.
\newcommand{\pgn}[1]{{\scriptsize\textbf{[#1]}}}

% Asterisco do original (divisão de assunto, sem título): \asterisco centrado em linha própria,
% também válido no meio do parágrafo ou dentro de nota (\footnote{... \asterisco ...}).
\newcommand{\asterisco}{\par\medskip\centerline{\footnotesize$\ast$}\medskip\noindent}

% Gravura recortada do fac-símile (<figure> do TEI; tools/tei2latex.py): fora de float, no ponto
% exato do original.  \gravura[detalhes]{arquivo.png}{fração da largura do texto}{legenda}:
% legenda principal centrada ("FIGURA VIII. — Côrte…"); os detalhes opcionais (<p> do <figure>:
% "A, hemispherio…; — B, …") vão abaixo, em corpo menor e justificados.  Ambos podem ser vazios.
\usepackage{graphicx}
\newcommand{\gravura}[4][]{\par\medskip\begingroup\centering
  \includegraphics[width=#3\linewidth,height=0.75\textheight,keepaspectratio]{#2}\par
  \if\relax\detokenize{#4}\relax\else\smallskip{\small #4\par}\fi
  \if\relax\detokenize{#1}\relax\else\smallskip{\footnotesize\leftskip=1em\rightskip=1em
    \parfillskip=0pt plus 1fil\noindent #1\par}\fi\endgroup\medskip\noindent}

% Citações latinas de S. Tomás (o OCR erra muito aqui: sempre conferir na imagem).
\newenvironment{latim}{\begin{quote}\itshape}{\end{quote}}

% --- Estrutura do autor (promovido de "Artigo III.tex") -----------------------
% enunciado de tese numerada que corre com o texto (segue-se travessão)
\newcommand{\tese}[2]{\bigskip\noindent\textbf{#1.\quad #2}\;\textemdash\;}
% enunciado de tese numerada em linha própria (seguem-se as subdivisões a), b)...)
\newcommand{\teseB}[2]{\bigskip\noindent\textbf{#1.\quad #2}\par}
% subdivisão a), b), c)... com rubrica em itálico
\newcommand{\sub}[2]{\medskip\noindent#1)\quad\emph{#2},\;\textemdash\;}
% subdivisão a), b), c)... sem rubrica
\newcommand{\subp}[1]{\medskip\noindent#1)\quad}

% --- Notas: no pé (padrão) ou correndo no corpo do texto ---------------------
% Passar  \def\notascorridas{}  no -usepretex faz cada nota correr no corpo, em
% corpo menor, num bloco recuado que interrompe a leitura até acabar.  Ex.:
%   latexmk -lualatex -usepretex='\def\fragmento{vol2/.../01_cap_y.tex}\def\notascorridas{}' tex/_preview.tex
\newif\ifnotascorridas
\ifdefined\notascorridas \notascorridastrue \fi
\ifnotascorridas
  \renewcommand\footnote[1]{%
    \stepcounter{footnote}\textsuperscript{\thefootnote}% marca no ponto do texto
    \par
    {\footnotesize
     \leftskip=2em \rightskip=2em \parindent=1em
     \addvspace{\medskipamount}%
     \noindent\textit{Nota \thefootnote.}\enspace #1\par
     \addvspace{\medskipamount}}%
  }
\fi

% --- Sem notas: \def\semnotas{} no -usepretex suprime todas as notas ---------
% (nem a chamada no texto nem o conteúdo).  Usado por  ./compila.sh --sem-notas .
\ifdefined\semnotas
  \renewcommand\footnote[2][]{}%
  \renewcommand\footnotemark[1][]{}%
  \renewcommand\footnotetext[2][]{}%
\fi

% peças pré-textuais (nó pretextual; p@type em web/obra/esquema.py)
\newcommand{\preTituloObra}[1]{\par\vspace{1em}{\centering\Large\bfseries\scshape #1\par}\vspace{1em}}
\newcommand{\preCentro}[1]{\par{\centering #1\par}}
\newcommand{\preEnderecamento}[1]{\par\vspace{1em}{\centering\small\scshape #1\par}}
\newcommand{\preSaudacao}[1]{\par\medskip\noindent #1\par}
\newcommand{\preDataLocal}[1]{\par\medskip\noindent\emph{#1}\par}
\newcommand{\preFecho}[1]{\par\smallskip{\raggedleft\scshape #1\par}}
}
\def\fimavulso{\end{document}}
\begin{document}
}{% Compilação AVULSA de um fragmento (volN/.../<stem>.tex) — p. ex. pelo botão do editor,
% de dentro da pasta do nó.  O fragmento começa com
%
%   \ifdefined\pg\else\IfFileExists{tex/avulso.tex}{% Compilação AVULSA de um fragmento (volN/.../<stem>.tex) — p. ex. pelo botão do editor,
% de dentro da pasta do nó.  O fragmento começa com
%
%   \ifdefined\pg\else\IfFileExists{tex/avulso.tex}{% Compilação AVULSA de um fragmento (volN/.../<stem>.tex) — p. ex. pelo botão do editor,
% de dentro da pasta do nó.  O fragmento começa com
%
%   \ifdefined\pg\else\IfFileExists{tex/avulso.tex}{\input{tex/avulso}}{\input{../../tex/avulso}}\fi
%
% e termina com  \csname fimavulso\endcsname .  Dentro do livro ou do _preview, \pg já existe
% e nada disto roda; \fimavulso fica indefinido e o \csname vira \relax.
%
% Motor: o projeto é lualatex (fontspec/polyglossia).  Se o editor chamar pdflatex, cai
% num preâmbulo equivalente com inputenc/babel — serve para conferir, não para o PDF final.
\documentclass[11pt,a5paper]{book}
\ifdefined\directlua
  \usepackage{fontspec}
  \usepackage{polyglossia}
  \setmainlanguage{portuguese}
\else\ifdefined\XeTeXrevision
  \usepackage{fontspec}
  \usepackage{polyglossia}
  \setmainlanguage{portuguese}
\else
  \usepackage[utf8]{inputenc}
  \usepackage[T1]{fontenc}
  \usepackage{lmodern}
  \usepackage[portuguese]{babel}
  \DeclareUnicodeCharacter{2192}{\ensuremath{\rightarrow}}
\fi\fi
% raiz do repo: compilado da raiz ou de volN/<parte>/ (dois níveis abaixo)
\IfFileExists{tex/preambulo.tex}{\input{tex/preambulo}}{\input{../../tex/preambulo}}
\def\fimavulso{\end{document}}
\begin{document}
}{% Compilação AVULSA de um fragmento (volN/.../<stem>.tex) — p. ex. pelo botão do editor,
% de dentro da pasta do nó.  O fragmento começa com
%
%   \ifdefined\pg\else\IfFileExists{tex/avulso.tex}{\input{tex/avulso}}{\input{../../tex/avulso}}\fi
%
% e termina com  \csname fimavulso\endcsname .  Dentro do livro ou do _preview, \pg já existe
% e nada disto roda; \fimavulso fica indefinido e o \csname vira \relax.
%
% Motor: o projeto é lualatex (fontspec/polyglossia).  Se o editor chamar pdflatex, cai
% num preâmbulo equivalente com inputenc/babel — serve para conferir, não para o PDF final.
\documentclass[11pt,a5paper]{book}
\ifdefined\directlua
  \usepackage{fontspec}
  \usepackage{polyglossia}
  \setmainlanguage{portuguese}
\else\ifdefined\XeTeXrevision
  \usepackage{fontspec}
  \usepackage{polyglossia}
  \setmainlanguage{portuguese}
\else
  \usepackage[utf8]{inputenc}
  \usepackage[T1]{fontenc}
  \usepackage{lmodern}
  \usepackage[portuguese]{babel}
  \DeclareUnicodeCharacter{2192}{\ensuremath{\rightarrow}}
\fi\fi
% raiz do repo: compilado da raiz ou de volN/<parte>/ (dois níveis abaixo)
\IfFileExists{tex/preambulo.tex}{\input{tex/preambulo}}{\input{../../tex/preambulo}}
\def\fimavulso{\end{document}}
\begin{document}
}\fi
%
% e termina com  \csname fimavulso\endcsname .  Dentro do livro ou do _preview, \pg já existe
% e nada disto roda; \fimavulso fica indefinido e o \csname vira \relax.
%
% Motor: o projeto é lualatex (fontspec/polyglossia).  Se o editor chamar pdflatex, cai
% num preâmbulo equivalente com inputenc/babel — serve para conferir, não para o PDF final.
\documentclass[11pt,a5paper]{book}
\ifdefined\directlua
  \usepackage{fontspec}
  \usepackage{polyglossia}
  \setmainlanguage{portuguese}
\else\ifdefined\XeTeXrevision
  \usepackage{fontspec}
  \usepackage{polyglossia}
  \setmainlanguage{portuguese}
\else
  \usepackage[utf8]{inputenc}
  \usepackage[T1]{fontenc}
  \usepackage{lmodern}
  \usepackage[portuguese]{babel}
  \DeclareUnicodeCharacter{2192}{\ensuremath{\rightarrow}}
\fi\fi
% raiz do repo: compilado da raiz ou de volN/<parte>/ (dois níveis abaixo)
\IfFileExists{tex/preambulo.tex}{% Preâmbulo comum ao livro e à pré-visualização de fragmentos (tex/_preview.tex).

% Mancha em A5.  Sem isto vale o padrão da classe book (~105 mm de largura e margem
% externa de ~25 mm em frente e verso), largo demais para o formato.  Aqui: ~118 mm de
% largura, margem interna um pouco maior que a externa (encadernação).
\usepackage[a5paper,inner=16mm,outer=14mm,top=14mm,bottom=16mm,
            includehead,headsep=5mm,footskip=8mm]{geometry}

% Títulos de secção sem justificar (evita linhas estouradas em \Large bold, p. ex.
% «Demonstrabilidade») e um pouco de folga para enunciados de tese em negrito.
\usepackage{sectsty}
\allsectionsfont{\raggedright}
\setlength{\emergencystretch}{1.5em}

% Marca da página impressa do original: \pg{286} imprime [286] no fluxo do texto.
\newcommand{\pg}[1]{\,{\footnotesize\textbf{[#1]}}\,}

% Virada de página do original DENTRO de uma nota (\pgn{422}); discreta, para não quebrar a nota.
\newcommand{\pgn}[1]{{\scriptsize\textbf{[#1]}}}

% Asterisco do original (divisão de assunto, sem título): \asterisco centrado em linha própria,
% também válido no meio do parágrafo ou dentro de nota (\footnote{... \asterisco ...}).
\newcommand{\asterisco}{\par\medskip\centerline{\footnotesize$\ast$}\medskip\noindent}

% Gravura recortada do fac-símile (<figure> do TEI; tools/tei2latex.py): fora de float, no ponto
% exato do original.  \gravura[detalhes]{arquivo.png}{fração da largura do texto}{legenda}:
% legenda principal centrada ("FIGURA VIII. — Côrte…"); os detalhes opcionais (<p> do <figure>:
% "A, hemispherio…; — B, …") vão abaixo, em corpo menor e justificados.  Ambos podem ser vazios.
\usepackage{graphicx}
\newcommand{\gravura}[4][]{\par\medskip\begingroup\centering
  \includegraphics[width=#3\linewidth,height=0.75\textheight,keepaspectratio]{#2}\par
  \if\relax\detokenize{#4}\relax\else\smallskip{\small #4\par}\fi
  \if\relax\detokenize{#1}\relax\else\smallskip{\footnotesize\leftskip=1em\rightskip=1em
    \parfillskip=0pt plus 1fil\noindent #1\par}\fi\endgroup\medskip\noindent}

% Citações latinas de S. Tomás (o OCR erra muito aqui: sempre conferir na imagem).
\newenvironment{latim}{\begin{quote}\itshape}{\end{quote}}

% --- Estrutura do autor (promovido de "Artigo III.tex") -----------------------
% enunciado de tese numerada que corre com o texto (segue-se travessão)
\newcommand{\tese}[2]{\bigskip\noindent\textbf{#1.\quad #2}\;\textemdash\;}
% enunciado de tese numerada em linha própria (seguem-se as subdivisões a), b)...)
\newcommand{\teseB}[2]{\bigskip\noindent\textbf{#1.\quad #2}\par}
% subdivisão a), b), c)... com rubrica em itálico
\newcommand{\sub}[2]{\medskip\noindent#1)\quad\emph{#2},\;\textemdash\;}
% subdivisão a), b), c)... sem rubrica
\newcommand{\subp}[1]{\medskip\noindent#1)\quad}

% --- Notas: no pé (padrão) ou correndo no corpo do texto ---------------------
% Passar  \def\notascorridas{}  no -usepretex faz cada nota correr no corpo, em
% corpo menor, num bloco recuado que interrompe a leitura até acabar.  Ex.:
%   latexmk -lualatex -usepretex='\def\fragmento{vol2/.../01_cap_y.tex}\def\notascorridas{}' tex/_preview.tex
\newif\ifnotascorridas
\ifdefined\notascorridas \notascorridastrue \fi
\ifnotascorridas
  \renewcommand\footnote[1]{%
    \stepcounter{footnote}\textsuperscript{\thefootnote}% marca no ponto do texto
    \par
    {\footnotesize
     \leftskip=2em \rightskip=2em \parindent=1em
     \addvspace{\medskipamount}%
     \noindent\textit{Nota \thefootnote.}\enspace #1\par
     \addvspace{\medskipamount}}%
  }
\fi

% --- Sem notas: \def\semnotas{} no -usepretex suprime todas as notas ---------
% (nem a chamada no texto nem o conteúdo).  Usado por  ./compila.sh --sem-notas .
\ifdefined\semnotas
  \renewcommand\footnote[2][]{}%
  \renewcommand\footnotemark[1][]{}%
  \renewcommand\footnotetext[2][]{}%
\fi

% peças pré-textuais (nó pretextual; p@type em web/obra/esquema.py)
\newcommand{\preTituloObra}[1]{\par\vspace{1em}{\centering\Large\bfseries\scshape #1\par}\vspace{1em}}
\newcommand{\preCentro}[1]{\par{\centering #1\par}}
\newcommand{\preEnderecamento}[1]{\par\vspace{1em}{\centering\small\scshape #1\par}}
\newcommand{\preSaudacao}[1]{\par\medskip\noindent #1\par}
\newcommand{\preDataLocal}[1]{\par\medskip\noindent\emph{#1}\par}
\newcommand{\preFecho}[1]{\par\smallskip{\raggedleft\scshape #1\par}}
}{% Preâmbulo comum ao livro e à pré-visualização de fragmentos (tex/_preview.tex).

% Mancha em A5.  Sem isto vale o padrão da classe book (~105 mm de largura e margem
% externa de ~25 mm em frente e verso), largo demais para o formato.  Aqui: ~118 mm de
% largura, margem interna um pouco maior que a externa (encadernação).
\usepackage[a5paper,inner=16mm,outer=14mm,top=14mm,bottom=16mm,
            includehead,headsep=5mm,footskip=8mm]{geometry}

% Títulos de secção sem justificar (evita linhas estouradas em \Large bold, p. ex.
% «Demonstrabilidade») e um pouco de folga para enunciados de tese em negrito.
\usepackage{sectsty}
\allsectionsfont{\raggedright}
\setlength{\emergencystretch}{1.5em}

% Marca da página impressa do original: \pg{286} imprime [286] no fluxo do texto.
\newcommand{\pg}[1]{\,{\footnotesize\textbf{[#1]}}\,}

% Virada de página do original DENTRO de uma nota (\pgn{422}); discreta, para não quebrar a nota.
\newcommand{\pgn}[1]{{\scriptsize\textbf{[#1]}}}

% Asterisco do original (divisão de assunto, sem título): \asterisco centrado em linha própria,
% também válido no meio do parágrafo ou dentro de nota (\footnote{... \asterisco ...}).
\newcommand{\asterisco}{\par\medskip\centerline{\footnotesize$\ast$}\medskip\noindent}

% Gravura recortada do fac-símile (<figure> do TEI; tools/tei2latex.py): fora de float, no ponto
% exato do original.  \gravura[detalhes]{arquivo.png}{fração da largura do texto}{legenda}:
% legenda principal centrada ("FIGURA VIII. — Côrte…"); os detalhes opcionais (<p> do <figure>:
% "A, hemispherio…; — B, …") vão abaixo, em corpo menor e justificados.  Ambos podem ser vazios.
\usepackage{graphicx}
\newcommand{\gravura}[4][]{\par\medskip\begingroup\centering
  \includegraphics[width=#3\linewidth,height=0.75\textheight,keepaspectratio]{#2}\par
  \if\relax\detokenize{#4}\relax\else\smallskip{\small #4\par}\fi
  \if\relax\detokenize{#1}\relax\else\smallskip{\footnotesize\leftskip=1em\rightskip=1em
    \parfillskip=0pt plus 1fil\noindent #1\par}\fi\endgroup\medskip\noindent}

% Citações latinas de S. Tomás (o OCR erra muito aqui: sempre conferir na imagem).
\newenvironment{latim}{\begin{quote}\itshape}{\end{quote}}

% --- Estrutura do autor (promovido de "Artigo III.tex") -----------------------
% enunciado de tese numerada que corre com o texto (segue-se travessão)
\newcommand{\tese}[2]{\bigskip\noindent\textbf{#1.\quad #2}\;\textemdash\;}
% enunciado de tese numerada em linha própria (seguem-se as subdivisões a), b)...)
\newcommand{\teseB}[2]{\bigskip\noindent\textbf{#1.\quad #2}\par}
% subdivisão a), b), c)... com rubrica em itálico
\newcommand{\sub}[2]{\medskip\noindent#1)\quad\emph{#2},\;\textemdash\;}
% subdivisão a), b), c)... sem rubrica
\newcommand{\subp}[1]{\medskip\noindent#1)\quad}

% --- Notas: no pé (padrão) ou correndo no corpo do texto ---------------------
% Passar  \def\notascorridas{}  no -usepretex faz cada nota correr no corpo, em
% corpo menor, num bloco recuado que interrompe a leitura até acabar.  Ex.:
%   latexmk -lualatex -usepretex='\def\fragmento{vol2/.../01_cap_y.tex}\def\notascorridas{}' tex/_preview.tex
\newif\ifnotascorridas
\ifdefined\notascorridas \notascorridastrue \fi
\ifnotascorridas
  \renewcommand\footnote[1]{%
    \stepcounter{footnote}\textsuperscript{\thefootnote}% marca no ponto do texto
    \par
    {\footnotesize
     \leftskip=2em \rightskip=2em \parindent=1em
     \addvspace{\medskipamount}%
     \noindent\textit{Nota \thefootnote.}\enspace #1\par
     \addvspace{\medskipamount}}%
  }
\fi

% --- Sem notas: \def\semnotas{} no -usepretex suprime todas as notas ---------
% (nem a chamada no texto nem o conteúdo).  Usado por  ./compila.sh --sem-notas .
\ifdefined\semnotas
  \renewcommand\footnote[2][]{}%
  \renewcommand\footnotemark[1][]{}%
  \renewcommand\footnotetext[2][]{}%
\fi

% peças pré-textuais (nó pretextual; p@type em web/obra/esquema.py)
\newcommand{\preTituloObra}[1]{\par\vspace{1em}{\centering\Large\bfseries\scshape #1\par}\vspace{1em}}
\newcommand{\preCentro}[1]{\par{\centering #1\par}}
\newcommand{\preEnderecamento}[1]{\par\vspace{1em}{\centering\small\scshape #1\par}}
\newcommand{\preSaudacao}[1]{\par\medskip\noindent #1\par}
\newcommand{\preDataLocal}[1]{\par\medskip\noindent\emph{#1}\par}
\newcommand{\preFecho}[1]{\par\smallskip{\raggedleft\scshape #1\par}}
}
\def\fimavulso{\end{document}}
\begin{document}
}{% Compilação AVULSA de um fragmento (volN/.../<stem>.tex) — p. ex. pelo botão do editor,
% de dentro da pasta do nó.  O fragmento começa com
%
%   \ifdefined\pg\else\IfFileExists{tex/avulso.tex}{% Compilação AVULSA de um fragmento (volN/.../<stem>.tex) — p. ex. pelo botão do editor,
% de dentro da pasta do nó.  O fragmento começa com
%
%   \ifdefined\pg\else\IfFileExists{tex/avulso.tex}{\input{tex/avulso}}{\input{../../tex/avulso}}\fi
%
% e termina com  \csname fimavulso\endcsname .  Dentro do livro ou do _preview, \pg já existe
% e nada disto roda; \fimavulso fica indefinido e o \csname vira \relax.
%
% Motor: o projeto é lualatex (fontspec/polyglossia).  Se o editor chamar pdflatex, cai
% num preâmbulo equivalente com inputenc/babel — serve para conferir, não para o PDF final.
\documentclass[11pt,a5paper]{book}
\ifdefined\directlua
  \usepackage{fontspec}
  \usepackage{polyglossia}
  \setmainlanguage{portuguese}
\else\ifdefined\XeTeXrevision
  \usepackage{fontspec}
  \usepackage{polyglossia}
  \setmainlanguage{portuguese}
\else
  \usepackage[utf8]{inputenc}
  \usepackage[T1]{fontenc}
  \usepackage{lmodern}
  \usepackage[portuguese]{babel}
  \DeclareUnicodeCharacter{2192}{\ensuremath{\rightarrow}}
\fi\fi
% raiz do repo: compilado da raiz ou de volN/<parte>/ (dois níveis abaixo)
\IfFileExists{tex/preambulo.tex}{\input{tex/preambulo}}{\input{../../tex/preambulo}}
\def\fimavulso{\end{document}}
\begin{document}
}{% Compilação AVULSA de um fragmento (volN/.../<stem>.tex) — p. ex. pelo botão do editor,
% de dentro da pasta do nó.  O fragmento começa com
%
%   \ifdefined\pg\else\IfFileExists{tex/avulso.tex}{\input{tex/avulso}}{\input{../../tex/avulso}}\fi
%
% e termina com  \csname fimavulso\endcsname .  Dentro do livro ou do _preview, \pg já existe
% e nada disto roda; \fimavulso fica indefinido e o \csname vira \relax.
%
% Motor: o projeto é lualatex (fontspec/polyglossia).  Se o editor chamar pdflatex, cai
% num preâmbulo equivalente com inputenc/babel — serve para conferir, não para o PDF final.
\documentclass[11pt,a5paper]{book}
\ifdefined\directlua
  \usepackage{fontspec}
  \usepackage{polyglossia}
  \setmainlanguage{portuguese}
\else\ifdefined\XeTeXrevision
  \usepackage{fontspec}
  \usepackage{polyglossia}
  \setmainlanguage{portuguese}
\else
  \usepackage[utf8]{inputenc}
  \usepackage[T1]{fontenc}
  \usepackage{lmodern}
  \usepackage[portuguese]{babel}
  \DeclareUnicodeCharacter{2192}{\ensuremath{\rightarrow}}
\fi\fi
% raiz do repo: compilado da raiz ou de volN/<parte>/ (dois níveis abaixo)
\IfFileExists{tex/preambulo.tex}{\input{tex/preambulo}}{\input{../../tex/preambulo}}
\def\fimavulso{\end{document}}
\begin{document}
}\fi
%
% e termina com  \csname fimavulso\endcsname .  Dentro do livro ou do _preview, \pg já existe
% e nada disto roda; \fimavulso fica indefinido e o \csname vira \relax.
%
% Motor: o projeto é lualatex (fontspec/polyglossia).  Se o editor chamar pdflatex, cai
% num preâmbulo equivalente com inputenc/babel — serve para conferir, não para o PDF final.
\documentclass[11pt,a5paper]{book}
\ifdefined\directlua
  \usepackage{fontspec}
  \usepackage{polyglossia}
  \setmainlanguage{portuguese}
\else\ifdefined\XeTeXrevision
  \usepackage{fontspec}
  \usepackage{polyglossia}
  \setmainlanguage{portuguese}
\else
  \usepackage[utf8]{inputenc}
  \usepackage[T1]{fontenc}
  \usepackage{lmodern}
  \usepackage[portuguese]{babel}
  \DeclareUnicodeCharacter{2192}{\ensuremath{\rightarrow}}
\fi\fi
% raiz do repo: compilado da raiz ou de volN/<parte>/ (dois níveis abaixo)
\IfFileExists{tex/preambulo.tex}{% Preâmbulo comum ao livro e à pré-visualização de fragmentos (tex/_preview.tex).

% Mancha em A5.  Sem isto vale o padrão da classe book (~105 mm de largura e margem
% externa de ~25 mm em frente e verso), largo demais para o formato.  Aqui: ~118 mm de
% largura, margem interna um pouco maior que a externa (encadernação).
\usepackage[a5paper,inner=16mm,outer=14mm,top=14mm,bottom=16mm,
            includehead,headsep=5mm,footskip=8mm]{geometry}

% Títulos de secção sem justificar (evita linhas estouradas em \Large bold, p. ex.
% «Demonstrabilidade») e um pouco de folga para enunciados de tese em negrito.
\usepackage{sectsty}
\allsectionsfont{\raggedright}
\setlength{\emergencystretch}{1.5em}

% Marca da página impressa do original: \pg{286} imprime [286] no fluxo do texto.
\newcommand{\pg}[1]{\,{\footnotesize\textbf{[#1]}}\,}

% Virada de página do original DENTRO de uma nota (\pgn{422}); discreta, para não quebrar a nota.
\newcommand{\pgn}[1]{{\scriptsize\textbf{[#1]}}}

% Asterisco do original (divisão de assunto, sem título): \asterisco centrado em linha própria,
% também válido no meio do parágrafo ou dentro de nota (\footnote{... \asterisco ...}).
\newcommand{\asterisco}{\par\medskip\centerline{\footnotesize$\ast$}\medskip\noindent}

% Gravura recortada do fac-símile (<figure> do TEI; tools/tei2latex.py): fora de float, no ponto
% exato do original.  \gravura[detalhes]{arquivo.png}{fração da largura do texto}{legenda}:
% legenda principal centrada ("FIGURA VIII. — Côrte…"); os detalhes opcionais (<p> do <figure>:
% "A, hemispherio…; — B, …") vão abaixo, em corpo menor e justificados.  Ambos podem ser vazios.
\usepackage{graphicx}
\newcommand{\gravura}[4][]{\par\medskip\begingroup\centering
  \includegraphics[width=#3\linewidth,height=0.75\textheight,keepaspectratio]{#2}\par
  \if\relax\detokenize{#4}\relax\else\smallskip{\small #4\par}\fi
  \if\relax\detokenize{#1}\relax\else\smallskip{\footnotesize\leftskip=1em\rightskip=1em
    \parfillskip=0pt plus 1fil\noindent #1\par}\fi\endgroup\medskip\noindent}

% Citações latinas de S. Tomás (o OCR erra muito aqui: sempre conferir na imagem).
\newenvironment{latim}{\begin{quote}\itshape}{\end{quote}}

% --- Estrutura do autor (promovido de "Artigo III.tex") -----------------------
% enunciado de tese numerada que corre com o texto (segue-se travessão)
\newcommand{\tese}[2]{\bigskip\noindent\textbf{#1.\quad #2}\;\textemdash\;}
% enunciado de tese numerada em linha própria (seguem-se as subdivisões a), b)...)
\newcommand{\teseB}[2]{\bigskip\noindent\textbf{#1.\quad #2}\par}
% subdivisão a), b), c)... com rubrica em itálico
\newcommand{\sub}[2]{\medskip\noindent#1)\quad\emph{#2},\;\textemdash\;}
% subdivisão a), b), c)... sem rubrica
\newcommand{\subp}[1]{\medskip\noindent#1)\quad}

% --- Notas: no pé (padrão) ou correndo no corpo do texto ---------------------
% Passar  \def\notascorridas{}  no -usepretex faz cada nota correr no corpo, em
% corpo menor, num bloco recuado que interrompe a leitura até acabar.  Ex.:
%   latexmk -lualatex -usepretex='\def\fragmento{vol2/.../01_cap_y.tex}\def\notascorridas{}' tex/_preview.tex
\newif\ifnotascorridas
\ifdefined\notascorridas \notascorridastrue \fi
\ifnotascorridas
  \renewcommand\footnote[1]{%
    \stepcounter{footnote}\textsuperscript{\thefootnote}% marca no ponto do texto
    \par
    {\footnotesize
     \leftskip=2em \rightskip=2em \parindent=1em
     \addvspace{\medskipamount}%
     \noindent\textit{Nota \thefootnote.}\enspace #1\par
     \addvspace{\medskipamount}}%
  }
\fi

% --- Sem notas: \def\semnotas{} no -usepretex suprime todas as notas ---------
% (nem a chamada no texto nem o conteúdo).  Usado por  ./compila.sh --sem-notas .
\ifdefined\semnotas
  \renewcommand\footnote[2][]{}%
  \renewcommand\footnotemark[1][]{}%
  \renewcommand\footnotetext[2][]{}%
\fi

% peças pré-textuais (nó pretextual; p@type em web/obra/esquema.py)
\newcommand{\preTituloObra}[1]{\par\vspace{1em}{\centering\Large\bfseries\scshape #1\par}\vspace{1em}}
\newcommand{\preCentro}[1]{\par{\centering #1\par}}
\newcommand{\preEnderecamento}[1]{\par\vspace{1em}{\centering\small\scshape #1\par}}
\newcommand{\preSaudacao}[1]{\par\medskip\noindent #1\par}
\newcommand{\preDataLocal}[1]{\par\medskip\noindent\emph{#1}\par}
\newcommand{\preFecho}[1]{\par\smallskip{\raggedleft\scshape #1\par}}
}{% Preâmbulo comum ao livro e à pré-visualização de fragmentos (tex/_preview.tex).

% Mancha em A5.  Sem isto vale o padrão da classe book (~105 mm de largura e margem
% externa de ~25 mm em frente e verso), largo demais para o formato.  Aqui: ~118 mm de
% largura, margem interna um pouco maior que a externa (encadernação).
\usepackage[a5paper,inner=16mm,outer=14mm,top=14mm,bottom=16mm,
            includehead,headsep=5mm,footskip=8mm]{geometry}

% Títulos de secção sem justificar (evita linhas estouradas em \Large bold, p. ex.
% «Demonstrabilidade») e um pouco de folga para enunciados de tese em negrito.
\usepackage{sectsty}
\allsectionsfont{\raggedright}
\setlength{\emergencystretch}{1.5em}

% Marca da página impressa do original: \pg{286} imprime [286] no fluxo do texto.
\newcommand{\pg}[1]{\,{\footnotesize\textbf{[#1]}}\,}

% Virada de página do original DENTRO de uma nota (\pgn{422}); discreta, para não quebrar a nota.
\newcommand{\pgn}[1]{{\scriptsize\textbf{[#1]}}}

% Asterisco do original (divisão de assunto, sem título): \asterisco centrado em linha própria,
% também válido no meio do parágrafo ou dentro de nota (\footnote{... \asterisco ...}).
\newcommand{\asterisco}{\par\medskip\centerline{\footnotesize$\ast$}\medskip\noindent}

% Gravura recortada do fac-símile (<figure> do TEI; tools/tei2latex.py): fora de float, no ponto
% exato do original.  \gravura[detalhes]{arquivo.png}{fração da largura do texto}{legenda}:
% legenda principal centrada ("FIGURA VIII. — Côrte…"); os detalhes opcionais (<p> do <figure>:
% "A, hemispherio…; — B, …") vão abaixo, em corpo menor e justificados.  Ambos podem ser vazios.
\usepackage{graphicx}
\newcommand{\gravura}[4][]{\par\medskip\begingroup\centering
  \includegraphics[width=#3\linewidth,height=0.75\textheight,keepaspectratio]{#2}\par
  \if\relax\detokenize{#4}\relax\else\smallskip{\small #4\par}\fi
  \if\relax\detokenize{#1}\relax\else\smallskip{\footnotesize\leftskip=1em\rightskip=1em
    \parfillskip=0pt plus 1fil\noindent #1\par}\fi\endgroup\medskip\noindent}

% Citações latinas de S. Tomás (o OCR erra muito aqui: sempre conferir na imagem).
\newenvironment{latim}{\begin{quote}\itshape}{\end{quote}}

% --- Estrutura do autor (promovido de "Artigo III.tex") -----------------------
% enunciado de tese numerada que corre com o texto (segue-se travessão)
\newcommand{\tese}[2]{\bigskip\noindent\textbf{#1.\quad #2}\;\textemdash\;}
% enunciado de tese numerada em linha própria (seguem-se as subdivisões a), b)...)
\newcommand{\teseB}[2]{\bigskip\noindent\textbf{#1.\quad #2}\par}
% subdivisão a), b), c)... com rubrica em itálico
\newcommand{\sub}[2]{\medskip\noindent#1)\quad\emph{#2},\;\textemdash\;}
% subdivisão a), b), c)... sem rubrica
\newcommand{\subp}[1]{\medskip\noindent#1)\quad}

% --- Notas: no pé (padrão) ou correndo no corpo do texto ---------------------
% Passar  \def\notascorridas{}  no -usepretex faz cada nota correr no corpo, em
% corpo menor, num bloco recuado que interrompe a leitura até acabar.  Ex.:
%   latexmk -lualatex -usepretex='\def\fragmento{vol2/.../01_cap_y.tex}\def\notascorridas{}' tex/_preview.tex
\newif\ifnotascorridas
\ifdefined\notascorridas \notascorridastrue \fi
\ifnotascorridas
  \renewcommand\footnote[1]{%
    \stepcounter{footnote}\textsuperscript{\thefootnote}% marca no ponto do texto
    \par
    {\footnotesize
     \leftskip=2em \rightskip=2em \parindent=1em
     \addvspace{\medskipamount}%
     \noindent\textit{Nota \thefootnote.}\enspace #1\par
     \addvspace{\medskipamount}}%
  }
\fi

% --- Sem notas: \def\semnotas{} no -usepretex suprime todas as notas ---------
% (nem a chamada no texto nem o conteúdo).  Usado por  ./compila.sh --sem-notas .
\ifdefined\semnotas
  \renewcommand\footnote[2][]{}%
  \renewcommand\footnotemark[1][]{}%
  \renewcommand\footnotetext[2][]{}%
\fi

% peças pré-textuais (nó pretextual; p@type em web/obra/esquema.py)
\newcommand{\preTituloObra}[1]{\par\vspace{1em}{\centering\Large\bfseries\scshape #1\par}\vspace{1em}}
\newcommand{\preCentro}[1]{\par{\centering #1\par}}
\newcommand{\preEnderecamento}[1]{\par\vspace{1em}{\centering\small\scshape #1\par}}
\newcommand{\preSaudacao}[1]{\par\medskip\noindent #1\par}
\newcommand{\preDataLocal}[1]{\par\medskip\noindent\emph{#1}\par}
\newcommand{\preFecho}[1]{\par\smallskip{\raggedleft\scshape #1\par}}
}
\def\fimavulso{\end{document}}
\begin{document}
}\fi
%
% e termina com  \csname fimavulso\endcsname .  Dentro do livro ou do _preview, \pg já existe
% e nada disto roda; \fimavulso fica indefinido e o \csname vira \relax.
%
% Motor: o projeto é lualatex (fontspec/polyglossia).  Se o editor chamar pdflatex, cai
% num preâmbulo equivalente com inputenc/babel — serve para conferir, não para o PDF final.
\documentclass[11pt,a5paper]{book}
\ifdefined\directlua
  \usepackage{fontspec}
  \usepackage{polyglossia}
  \setmainlanguage{portuguese}
\else\ifdefined\XeTeXrevision
  \usepackage{fontspec}
  \usepackage{polyglossia}
  \setmainlanguage{portuguese}
\else
  \usepackage[utf8]{inputenc}
  \usepackage[T1]{fontenc}
  \usepackage{lmodern}
  \usepackage[portuguese]{babel}
  \DeclareUnicodeCharacter{2192}{\ensuremath{\rightarrow}}
\fi\fi
% raiz do repo: compilado da raiz ou de volN/<parte>/ (dois níveis abaixo)
\IfFileExists{tex/preambulo.tex}{% Preâmbulo comum ao livro e à pré-visualização de fragmentos (tex/_preview.tex).

% Mancha em A5.  Sem isto vale o padrão da classe book (~105 mm de largura e margem
% externa de ~25 mm em frente e verso), largo demais para o formato.  Aqui: ~118 mm de
% largura, margem interna um pouco maior que a externa (encadernação).
\usepackage[a5paper,inner=16mm,outer=14mm,top=14mm,bottom=16mm,
            includehead,headsep=5mm,footskip=8mm]{geometry}

% Títulos de secção sem justificar (evita linhas estouradas em \Large bold, p. ex.
% «Demonstrabilidade») e um pouco de folga para enunciados de tese em negrito.
\usepackage{sectsty}
\allsectionsfont{\raggedright}
\setlength{\emergencystretch}{1.5em}

% Marca da página impressa do original: \pg{286} imprime [286] no fluxo do texto.
\newcommand{\pg}[1]{\,{\footnotesize\textbf{[#1]}}\,}

% Virada de página do original DENTRO de uma nota (\pgn{422}); discreta, para não quebrar a nota.
\newcommand{\pgn}[1]{{\scriptsize\textbf{[#1]}}}

% Asterisco do original (divisão de assunto, sem título): \asterisco centrado em linha própria,
% também válido no meio do parágrafo ou dentro de nota (\footnote{... \asterisco ...}).
\newcommand{\asterisco}{\par\medskip\centerline{\footnotesize$\ast$}\medskip\noindent}

% Gravura recortada do fac-símile (<figure> do TEI; tools/tei2latex.py): fora de float, no ponto
% exato do original.  \gravura[detalhes]{arquivo.png}{fração da largura do texto}{legenda}:
% legenda principal centrada ("FIGURA VIII. — Côrte…"); os detalhes opcionais (<p> do <figure>:
% "A, hemispherio…; — B, …") vão abaixo, em corpo menor e justificados.  Ambos podem ser vazios.
\usepackage{graphicx}
\newcommand{\gravura}[4][]{\par\medskip\begingroup\centering
  \includegraphics[width=#3\linewidth,height=0.75\textheight,keepaspectratio]{#2}\par
  \if\relax\detokenize{#4}\relax\else\smallskip{\small #4\par}\fi
  \if\relax\detokenize{#1}\relax\else\smallskip{\footnotesize\leftskip=1em\rightskip=1em
    \parfillskip=0pt plus 1fil\noindent #1\par}\fi\endgroup\medskip\noindent}

% Citações latinas de S. Tomás (o OCR erra muito aqui: sempre conferir na imagem).
\newenvironment{latim}{\begin{quote}\itshape}{\end{quote}}

% --- Estrutura do autor (promovido de "Artigo III.tex") -----------------------
% enunciado de tese numerada que corre com o texto (segue-se travessão)
\newcommand{\tese}[2]{\bigskip\noindent\textbf{#1.\quad #2}\;\textemdash\;}
% enunciado de tese numerada em linha própria (seguem-se as subdivisões a), b)...)
\newcommand{\teseB}[2]{\bigskip\noindent\textbf{#1.\quad #2}\par}
% subdivisão a), b), c)... com rubrica em itálico
\newcommand{\sub}[2]{\medskip\noindent#1)\quad\emph{#2},\;\textemdash\;}
% subdivisão a), b), c)... sem rubrica
\newcommand{\subp}[1]{\medskip\noindent#1)\quad}

% --- Notas: no pé (padrão) ou correndo no corpo do texto ---------------------
% Passar  \def\notascorridas{}  no -usepretex faz cada nota correr no corpo, em
% corpo menor, num bloco recuado que interrompe a leitura até acabar.  Ex.:
%   latexmk -lualatex -usepretex='\def\fragmento{vol2/.../01_cap_y.tex}\def\notascorridas{}' tex/_preview.tex
\newif\ifnotascorridas
\ifdefined\notascorridas \notascorridastrue \fi
\ifnotascorridas
  \renewcommand\footnote[1]{%
    \stepcounter{footnote}\textsuperscript{\thefootnote}% marca no ponto do texto
    \par
    {\footnotesize
     \leftskip=2em \rightskip=2em \parindent=1em
     \addvspace{\medskipamount}%
     \noindent\textit{Nota \thefootnote.}\enspace #1\par
     \addvspace{\medskipamount}}%
  }
\fi

% --- Sem notas: \def\semnotas{} no -usepretex suprime todas as notas ---------
% (nem a chamada no texto nem o conteúdo).  Usado por  ./compila.sh --sem-notas .
\ifdefined\semnotas
  \renewcommand\footnote[2][]{}%
  \renewcommand\footnotemark[1][]{}%
  \renewcommand\footnotetext[2][]{}%
\fi

% peças pré-textuais (nó pretextual; p@type em web/obra/esquema.py)
\newcommand{\preTituloObra}[1]{\par\vspace{1em}{\centering\Large\bfseries\scshape #1\par}\vspace{1em}}
\newcommand{\preCentro}[1]{\par{\centering #1\par}}
\newcommand{\preEnderecamento}[1]{\par\vspace{1em}{\centering\small\scshape #1\par}}
\newcommand{\preSaudacao}[1]{\par\medskip\noindent #1\par}
\newcommand{\preDataLocal}[1]{\par\medskip\noindent\emph{#1}\par}
\newcommand{\preFecho}[1]{\par\smallskip{\raggedleft\scshape #1\par}}
}{% Preâmbulo comum ao livro e à pré-visualização de fragmentos (tex/_preview.tex).

% Mancha em A5.  Sem isto vale o padrão da classe book (~105 mm de largura e margem
% externa de ~25 mm em frente e verso), largo demais para o formato.  Aqui: ~118 mm de
% largura, margem interna um pouco maior que a externa (encadernação).
\usepackage[a5paper,inner=16mm,outer=14mm,top=14mm,bottom=16mm,
            includehead,headsep=5mm,footskip=8mm]{geometry}

% Títulos de secção sem justificar (evita linhas estouradas em \Large bold, p. ex.
% «Demonstrabilidade») e um pouco de folga para enunciados de tese em negrito.
\usepackage{sectsty}
\allsectionsfont{\raggedright}
\setlength{\emergencystretch}{1.5em}

% Marca da página impressa do original: \pg{286} imprime [286] no fluxo do texto.
\newcommand{\pg}[1]{\,{\footnotesize\textbf{[#1]}}\,}

% Virada de página do original DENTRO de uma nota (\pgn{422}); discreta, para não quebrar a nota.
\newcommand{\pgn}[1]{{\scriptsize\textbf{[#1]}}}

% Asterisco do original (divisão de assunto, sem título): \asterisco centrado em linha própria,
% também válido no meio do parágrafo ou dentro de nota (\footnote{... \asterisco ...}).
\newcommand{\asterisco}{\par\medskip\centerline{\footnotesize$\ast$}\medskip\noindent}

% Gravura recortada do fac-símile (<figure> do TEI; tools/tei2latex.py): fora de float, no ponto
% exato do original.  \gravura[detalhes]{arquivo.png}{fração da largura do texto}{legenda}:
% legenda principal centrada ("FIGURA VIII. — Côrte…"); os detalhes opcionais (<p> do <figure>:
% "A, hemispherio…; — B, …") vão abaixo, em corpo menor e justificados.  Ambos podem ser vazios.
\usepackage{graphicx}
\newcommand{\gravura}[4][]{\par\medskip\begingroup\centering
  \includegraphics[width=#3\linewidth,height=0.75\textheight,keepaspectratio]{#2}\par
  \if\relax\detokenize{#4}\relax\else\smallskip{\small #4\par}\fi
  \if\relax\detokenize{#1}\relax\else\smallskip{\footnotesize\leftskip=1em\rightskip=1em
    \parfillskip=0pt plus 1fil\noindent #1\par}\fi\endgroup\medskip\noindent}

% Citações latinas de S. Tomás (o OCR erra muito aqui: sempre conferir na imagem).
\newenvironment{latim}{\begin{quote}\itshape}{\end{quote}}

% --- Estrutura do autor (promovido de "Artigo III.tex") -----------------------
% enunciado de tese numerada que corre com o texto (segue-se travessão)
\newcommand{\tese}[2]{\bigskip\noindent\textbf{#1.\quad #2}\;\textemdash\;}
% enunciado de tese numerada em linha própria (seguem-se as subdivisões a), b)...)
\newcommand{\teseB}[2]{\bigskip\noindent\textbf{#1.\quad #2}\par}
% subdivisão a), b), c)... com rubrica em itálico
\newcommand{\sub}[2]{\medskip\noindent#1)\quad\emph{#2},\;\textemdash\;}
% subdivisão a), b), c)... sem rubrica
\newcommand{\subp}[1]{\medskip\noindent#1)\quad}

% --- Notas: no pé (padrão) ou correndo no corpo do texto ---------------------
% Passar  \def\notascorridas{}  no -usepretex faz cada nota correr no corpo, em
% corpo menor, num bloco recuado que interrompe a leitura até acabar.  Ex.:
%   latexmk -lualatex -usepretex='\def\fragmento{vol2/.../01_cap_y.tex}\def\notascorridas{}' tex/_preview.tex
\newif\ifnotascorridas
\ifdefined\notascorridas \notascorridastrue \fi
\ifnotascorridas
  \renewcommand\footnote[1]{%
    \stepcounter{footnote}\textsuperscript{\thefootnote}% marca no ponto do texto
    \par
    {\footnotesize
     \leftskip=2em \rightskip=2em \parindent=1em
     \addvspace{\medskipamount}%
     \noindent\textit{Nota \thefootnote.}\enspace #1\par
     \addvspace{\medskipamount}}%
  }
\fi

% --- Sem notas: \def\semnotas{} no -usepretex suprime todas as notas ---------
% (nem a chamada no texto nem o conteúdo).  Usado por  ./compila.sh --sem-notas .
\ifdefined\semnotas
  \renewcommand\footnote[2][]{}%
  \renewcommand\footnotemark[1][]{}%
  \renewcommand\footnotetext[2][]{}%
\fi

% peças pré-textuais (nó pretextual; p@type em web/obra/esquema.py)
\newcommand{\preTituloObra}[1]{\par\vspace{1em}{\centering\Large\bfseries\scshape #1\par}\vspace{1em}}
\newcommand{\preCentro}[1]{\par{\centering #1\par}}
\newcommand{\preEnderecamento}[1]{\par\vspace{1em}{\centering\small\scshape #1\par}}
\newcommand{\preSaudacao}[1]{\par\medskip\noindent #1\par}
\newcommand{\preDataLocal}[1]{\par\medskip\noindent\emph{#1}\par}
\newcommand{\preFecho}[1]{\par\smallskip{\raggedleft\scshape #1\par}}
}
\def\fimavulso{\end{document}}
\begin{document}
}\fi
% PADRÃO-OURO vol. 1 (5.ª ed., 1927), bloco B — Metaphysica, Parte Primeira: Ontologia,
% Cap. II (Divisão do ente), Art. I.  Págs. impressas 331–334 (vol1/paginas/0331.png … 0334.png).
% Transcrito à mão da imagem; o OCR serviu só de rascunho.  Convenções em
% 062-065_logica_raciocinio.tex.  Aqui:
%   - o alto da p. 331 é o fim de uma nota chamada na p. 330: fica de fora;
%   - a nota (1) da p. 331 ocupa o pé das pp. 331–333; a nota (3) da p. 334 acaba na p. 335;
%   - « Deus » nas notas vem em negrito (\textbf), não em versalete;
%   - preservados: « Este infinito » (Summario), « sería ».

\pg{331}

\chapter*{Capitulo Segundo\\[0.3em]\large Divisão do ente}

{\small\noindent\textsc{Summario}: --- Ente necessario e ente contingente. --- Este infinito
e ente finito. --- Ente simples e ente composto. --- Ente immutavel e ente mutavel.\par}

\section*{Artigo I\\ Ente necessario e ente contingente.}

\tese{54}{Divisão do ente.} Todo o ente --- ou é \emph{acto puro}, --- ou consta de
\emph{acto} e de \emph{potencia}. Se é \emph{acto puro}, o ente é \emph{necessario,
infinito, simples, immutavel}. Se consta de \emph{acto} e de \emph{potencia}, é
\emph{contingente, finito, composto, mutavel}. D'onde a divisão de todo o ente --- em
\emph{necessario} e \emph{contingente}, --- em \emph{infinito} e \emph{finito}, --- em
\emph{simples} e \emph{composto}, --- em \emph{immutavel} e \emph{mutavel}. --- Esta
divisão funda-se na diversa perfeição do ente\footnote{Todo o ente é dotado de
\emph{perfeição}. --- \emph{Perfeição é toda a realidade, que serve} para fazer ou
completar uma coisa... A perfeição pode ser \emph{parcial} e \emph{total}. É
\emph{parcial}, se representa \emph{uma} das realidades, de que o ente é dotado. É
\emph{total}, se representa o \emph{conjuncto} das mesmas realidades. D'onde se vê que a
perfeição não é um attributo especial do ente, mas é uma coisa que constitue, adequada ou
inadequadamente, o proprio ente. --- Por isso um ente é e diz-se \emph{perfeito}, isto é,
\emph{inteiramente feito}, emquanto possue a perfeição \emph{total}, de modo que não lhe
falte nenhuma das realidades, necessarias para o seu complemento. --- A' \emph{perfeição}
oppõe-se a \emph{imperfeição}, que denota \emph{ausencia} de perfeição. A imperfeição é
\emph{privativa} ou \emph{negativa}. É \emph{privativa}, quando é ausencia de uma
perfeição \emph{devida}, e diz-se commummente \emph{privação}; tal é a ausencia de um
braço no homem. É \emph{negativa}, quando é a ausencia de uma perfeição
\emph{não-devida}, e chama-se \emph{negação}; tal é a ausencia das azas no homem, ou da
razão na planta.

A perfeição \emph{total} divide-se em \emph{absoluta} e \emph{relativa}. A
\emph{absoluta} é o conjuncto de \emph{todas} as realidades ou perfeições
\emph{possiveis}, sem limitação alguma; tal é a perfeição de \textbf{Deus}, que é o
proprio Ser subsistente. A \emph{relativa} é o conjuncto das realidades ou perfeições,
que são devidas a um certo genero, ou a uma certa especie de entes; tal é a perfeição de
um homem, que tem tudo o que a sua natureza exige. --- A perfeição \emph{absoluta} não
admitte graus, nem divisões; a \emph{relativa} admitte graus, --- emquanto um genero ou
uma especie pode abranger um maior ou menor numero de perfeições parciaes com relação a
um outro genero, ou a uma outra especie, --- ou emquanto um individuo possue uma
realidade ou perfeição por um modo mais \emph{intenso} que o modo por que a possue um
outro individuo do mesmo genero, ou da mesma especie. --- A perfeição \emph{total} pode
considerar-se --- em relação á \emph{entidade}, quando o ente tem tudo o que é
indispensavel para a sua essencia e integridade; --- em relação á \emph{operação},
quando o ente tem tudo o que é necessario para operar em conformidade com a sua
essencia; --- em relação ao \emph{fim}, quando o ente, tendo alcançado o fim pelas suas
operações, encontra nelle o seu complemento, pois o fim é bem (Cf. \emph{Sum. Th.}, p.
1, q. 6, a. 3).

A perfeição \emph{relativa} divide-se em \emph{simples} e \emph{mista}. ---
\emph{Simples} é a que no seu conceito não encerra imperfeição alguma e não se oppõe a
perfeições mais elevadas; tal é a \emph{vida}, a \emph{bondade}, a \emph{belleza}, etc. A
perfeição \emph{simples} costuma chamar-se tambem \emph{absoluta}, mas emquanto o seu
conceito não denota imperfeição. --- \emph{Mista} é a que, embora no seu conceito exprima
uma entidade positiva, todavia encerra alguma imperfeição e se oppõe a uma perfeição mais
elevada; tal é a \emph{corporeidade}, que, embora denote uma realidade, comtudo importa
uma imperfeição, pois o \emph{corpo} é \emph{composto, dependente}, etc., --- e se oppõe
á \emph{espiritualidade}, que é uma coisa mais elevada do que o corpo (Cf. Ansel., Monol.,
c. 15).

A perfeição pode convir ao ente por tres modos: \emph{formal, virtual} e
\emph{eminente}. Convem \emph{formalmente}, quando existe no ente segundo a propria forma
ou razão especifica, isto é, segundo o seu conceito e a sua definição; assim encontra-se
o \emph{calor no fogo}, a \emph{sensibilidade} no \emph{cão}, a \emph{razão} no
\emph{homem}. Convem \emph{virtualmente}, quando não existe no ente segundo a propria
forma, mas existe na virtude ou causalidade do ente, emquanto este é capaz de a produzir;
assim encontra-se a \emph{planta} na \emph{semente}, e todo o \emph{effeito} na sua
\emph{causa}. Convem \emph{eminentemente}; quando existe no ente, nem segundo a propria
forma, nem segundo a energia ou causalidade, mas por um modo mais nobre e mais elevado
que o modo por que se encontra nos seres, que a possuem \emph{formalmente}; assim
encontram-se em \textbf{Deus} todas as perfeições dos entes creados. --- Note-se que um
d'esses modos não exclue outro. Uma perfeição pode convir a um ente por um modo
\emph{formal} e \emph{eminente} ao mesmo tempo, e convem quando existe nelle segundo a
propria forma, mas num grau mais elevado; assim a \emph{sabedoria}, que é o conhecimento
das coisas pelas causas supremas, attribue-se a \textbf{Deus} --- não só
\emph{formalmente}, emquanto, em \textbf{Deus}, como no homem, importa o conhecimento
pelas causas supremas, --- mas tambem \emph{eminentemente}, emquanto em \textbf{Deus} a
\emph{sabedoria} é uma e a mesma coisa com essencia, ao passo que no homem é uma
qualidade accidental e limitada. Como tambem uma perfeição pode convir a um ente por um
modo \emph{virtual} e \emph{eminente} ao mesmo tempo, e convem quando não existe nelle
segundo a propria forma, mas existe na virtude ou capacidade do ente, e de um modo mais
elevado e nobre; assim o \emph{raciocinio}, importando um discurso ou uma passagem da
potencia para o acto, não se attribue a \textbf{Deus} \emph{formalmente}, --- mas
attribue-se \emph{virtualmente}, emquanto a intelligencia divina tem a capacidade de
perceber a conclusão no principio, --- e tambem \emph{eminentemente}, emquanto a
intelligencia divina não passa, como a do homem, do conhecido para o desconhecido, mas
vê a conclusão no principio sem discurso, por intuição.

Premissas estas noções, vejamos a razão da divisão do ente. A perfeição pode
encontrar-se no ente --- ou de um modo \emph{necessario} ou \emph{contingente}, e temos o
ente \emph{necessario} e o ente \emph{contingente}, --- ou de um modo \emph{illimitado} ou
\emph{limitado}, e temos o ente \emph{infinito} e o ente \emph{finito}, --- ou de um modo
\emph{simples} ou \emph{composto}, e temos o ente \emph{simples} e o ente
\emph{composto}, --- ou de um modo \emph{permanente} ou \emph{variavel}, e temos o ente
\emph{immutavel} e o ente \emph{mutavel}. --- A primeira especie da divisão refere-se á
\emph{existencia} do ente; a segunda, á sua \emph{comprehensão}; a terceira, á sua
\emph{constituição}; a quarta, á \emph{permanencia} ou \emph{estabilidade} da perfeição.
--- Como se vê, todas estas especies denotam a mesma coisa e importam os mesmos extremos,
considerados sob differentes aspectos e significados por varios conceitos e termos, dos
quaes um é mais claro e mais obvio que outro: porque o Ente \emph{necessario, infinito,
simples, immutavel} é \textbf{Deus}; e o ente \emph{contingente, finito, composto,
mutavel} é a creatura. Por isso esta divisão é a mais transcendental que se pode dar. ---
Advertimos que, tratando aqui do Ente \emph{necessario, infinito, simples} e
\emph{immutavel}, não entendemos demonstrar a sua \emph{existencia}; mas, admittida, pela
experiencia, a existencia de entes \emph{contingentes, finitos, compostos} e
\emph{mutaveis}, procuramos descobrir o \emph{nexo}, que elles têm com o Ente
\emph{necessario, infinito, simples e immutavel}.}.

\pg{332}
\tese{55}{Ente necessario e ente contingente.} Ente \emph{necessario} é o que existe
\emph{por si, pela exigencia da sua essencia}, e por isso não pode deixar de existir. ---
Ente \emph{contingente} é o que existe \emph{accidentalmente}, não pela exigencia da sua
essencia, mas pela causalidade de outro ente, e por isso é indifferente para a
\pg{333} existencia e para a não-existencia. --- No ente \emph{necessario}, a
\emph{existencia} é um elemento da \emph{essencia}, e identifica-se com esta; no ente
\emph{contingente}, a \emph{existencia} é \emph{indifferente} para a \emph{essencia}, e
uma coisa é \emph{realmente} distincta de outra\footnote{\emph{Ente necessario} é o que
não pode deixar de ser. Na verdade, \emph{ente} é o que possue o \emph{ser}, é o que tem
a \emph{existencia}, é o \emph{existente}. Portanto, o ente \emph{necessario} é o que
existe \emph{necessariamente}, e por isso possue o \emph{ser} por tal modo, que lhe
repugna o \emph{não-ser}. --- O \emph{ente contingente}, por ser \emph{contingente},
existe por tal modo, que pode deixar de existir, sem que d'ahi resulte algum absurdo ou
inconveniente. --- O ente \emph{necessario} pode ser \emph{absoluto} e
\emph{hypothetico}. É \emph{absoluto}, quando o \emph{não-ser} lhe repugna pela sua
intrinseca natureza, independentemente de toda a condição ou hypothese; tal é
\textbf{Deus}. É \emph{hypothetico}, quando o \emph{não-ser} lhe repugna, não pela sua
intrinseca natureza, mas debaixo de uma certa condição ou hypothese; tal é toda a
creatura, emquanto existe. --- Aqui fallamos do ente \emph{absoluto}.

O ente \emph{absoluto}, não tendo nenhuma potencia ou indifferença para a existencia, mas
existindo em virtude da propria essencia, é ente \emph{por si, pela sua essencia}, de
modo que a sua essencia é a razão sufficiente da sua existencia. --- Note-se que quando
se diz que o ente \emph{absoluto} existe pela \emph{sua essencia}, não se entende que a
essencia d'esse ente é a \emph{causa} da sua existencia: o que é absurdo; mas deve
entender-se que esse ente não depende, na sua existencia, de uma causa extrinseca e
superior. --- Advirta-se tambem que, embora a \emph{essencia} do ente \emph{absoluto} se
identifique com a \emph{existencia}, todavia a nossa intelligencia não percebe essas duas
coisas por uma só idea, mas por duas ideas distinctas, uma das quaes representa a
\emph{essencia}, e a outra a \emph{existencia}. E nada mais natural. A intelligencia
humana não tem um conhecimento intuitivo do Ente \emph{absoluto}, mas eleva-se a Elle por
meio das coisas creadas, e nestas a \emph{essencia é realmente} distincta da
\emph{existencia}.}.

\pg{334}
\teseB{56}{O ente, que é acto puro, é necessario; e o ente, que é necessario, é acto
puro.}

\sub{a}{O ente, que é acto puro, é necessario.} O ente, que é \emph{acto puro}, exclue de
si toda a potencia: aliás não sería o que é. Ora o ente, que exclue de si toda a
potencia, deve possuir a existencia identificada com a essencia, e por isso é
\emph{necessario}. Logo o ente, que é \emph{acto puro}, é \emph{necessario}\footnote{O
\emph{acto puro} exclue de si toda a potencia. Ora o ente, que exclue de si toda a
potencia, não pode ser \emph{contingente}; porque o ente \emph{contingente} não exclue,
mas encerra a potencia, emquanto está em potencia para o acto da existencia. Portanto o
ente, que é \emph{acto puro}, é \emph{necessario}, possue a existencia pela exigencia da
sua essencia. Nem podia ser diversamente. O ente, que é \emph{acto puro}, possue a
existencia identificada com a essencia, de modo que podemos dizer que a existencia é um
dos elementos da essencia. Sendo assim, esse ente é \emph{necessariamente existente},
como o homem, que pela sua essencia é dotado de razão, é \emph{necessariamente racional}.
As essencias, como vimos, são \emph{necessarias} quanto aos elementos, que as
constituem.}.

\sub{b}{O ente, que é necessario, é acto puro.} O ente, que é \emph{necessario}, existe
pela exigencia da sua essencia, e por isso possue a existencia identificada com a
essencia. Ora o ente, que possue a existencia identificada com a essencia, é \emph{acto
puro}. Logo o ente, que é \emph{necessario}, é \emph{acto puro}\footnote{O ente
\emph{necessario} é \emph{acto puro}. O \emph{acto} é \emph{puro}, quando não é recebido
num sujeito distincto, nem está em potencia para um acto ulterior. Ora o ente
\emph{necessario} --- não é recebido num sujeito distincto, porque o seu \emph{ser} é uma
e a mesma coisa com a \emph{essencia}, --- nem está em potencia para um acto ulterior,
porque o ente, cujo ser se identifica com a essencia, é o proprio ser, é o ser
subsistente, e por isso possue a plenitude da perfeição.

D'onde se vê que o ente \emph{necessario} é \emph{unico}. Por quanto o \emph{acto}, como
dissemos, só se multiplica emquanto é recebido em muitos sujeitos. Ora o ser do ente
\emph{necessario}, identificando-se com a essencia, não é recebido em nenhum sujeito, mas
é por si subsistente, e por isso não se multiplica, mas é \emph{unico}. Logo o ente
\emph{necessario} é \emph{unico}.}.

\teseB{57}{O ente, que consta de acto e de potencia, é contingente; e o ente, que é
contingente, consta de acto e de potencia.}

\sub{a}{O ente, que consta de acto e de potencia, é contingente.} O ente, que consta de
\emph{acto} e de \emph{potencia}, possue a existencia \emph{realmente} distincta da
essencia. Ora o ente, que possue a existencia \emph{realmente} distincta da essencia, é
\emph{contingente}. Logo o ente, que consta de \emph{acto} e de \emph{potencia}, é
\emph{contingente}\footnote{A conclusão é certissima. Na verdade, todo o ente composto
pode ser dissolvido; e, podendo ser dissolvido, é naturalmente \emph{contingente}. ---
Além d'isso, o ente composto é \emph{posterior} ás partes componentes, e por isso não pode
ser \emph{necessario}, mas deve ser \emph{contingente}; porque o ente \emph{necessario},
sendo o proprio ser, não está ordenado para outro ser mais perfeito, e é simplesmente o
\emph{primeiro}. --- Mais: o composto depende, no seu ser, de uma causa, que une as parte
componentes e que deve ser não só anterior ao proprio composto, mas distincta delle.
Dependendo, no seu ser, de uma causa distincta, o composto não existe pela sua natureza,
e por isso é \emph{contingente}.}.

\csname fimavulso\endcsname
